\chapter{فضاءات \texorpdfstring{$L^p$}{Lp}}\label{ch:b3:lp}

بُني تكامل لوبيغ من أجل التحليل؛ وفضاءات $L^p$ هي المكان الذي
يسكنه ذلك التحليل. فهي فضاءات باناخ (ريس--فيشر) --- أي
الإتمامات التي بيّن \cref{exo:b3:complete:1} أن الدوال المتصلة تفتقر إليها ---
وهي تحمل تقنية تنعيم، هي \emph{الالتفاف مع المليِّفات}، تقرّب
كل دالة من $L^p$ بدوال $\mathcal C^\infty$. ويبرهن هذا الفصل على صيغتي هولدر
ومينكوفسكي التكامليتين، وعلى التمام، ومبرهنات الكثافة، وآلة
التنظيم، وينتهي بجغرافيا الاحتواء والاستيفاء في سلّم $L^p$. وفي
كل ما يلي، يكون $(X, \mathcal A, \mu)$ فضاءً قياسيًّا وتكون الدوال ذات قيم عقدية؛
وعلى $\R^d$ يكون \omterm{def:b3:measure:measure}{القياس} هو $\lambda_d$.

\section{التعريف؛ هولدر ومينكوفسكي}

\begin{definition}\label{def:b3:lp:space}
من أجل $1 \leq p < \infty$، تكون $\mathcal L^p(\mu)$ مجموعة الدوال القابلة \omterm{def:b3:measure:measure}{للقياس} $f$ التي
تحقق $\norm f_p = \bigl(\int\abs
f^p\dd\mu\bigr)^{1/p} < \infty$، وتكون $\mathcal L^\infty(\mu)$ مجموعة الدوال $f$ المحدودة خارج مجموعة
معدومة \omterm{def:b3:measure:measure}{القياس}، مع $\norm
f_\infty$ \emph{السوپريموم الجوهري} --- أي أصغر
عدد $M$ يحقق $\abs f \leq M$ في كل مكان تقريبًا (والإنفيموم
مبلوغ: يكفي مقاطعة المجموعات المعدومة الموافقة للأعداد
$M + \frac1n$). وبما أن $\norm f_p = 0$ لا يفرض إلا $f
= 0$
\emph{في كل مكان تقريبًا} (\cref{exo:b3:lebesgue:5})، نعرّف
\[
L^p(\mu) = \mathcal L^p(\mu)/\{f = 0 \text{ في كل مكان تقريبًا}\} :
\]
فتكون العناصر أصنافَ دوال بترديد المجموعات المعدومة، ويكون
$\norm\cdot_p$ معيارًا فعليًّا على
$L^p$.\index{فضاء Lp@فضاء $L^p$}
\end{definition}

\begin{theorem}[متراجحة هولدر]\label{thm:b3:lp:holder}
ليكن $1 \leq p, q \leq \infty$ مع $\frac1p + \frac1q = 1$
(أي أُسّين مترافقين). من أجل دالتين قابلتين \omterm{def:b3:measure:measure}{للقياس} $f, g$:\index{متراجحة
هولدر}
\[
\norm{fg}_1 \leq \norm f_p\,\norm g_q ,
\]
مع التساوي (من أجل $1 < p < \infty$، ونظائم منتهية،
و $f,g\ne0$) إذا وفقط إذا كان $\abs f^p$ و $\abs g^q$ متناسبين
في كل مكان تقريبًا.
\end{theorem}

\begin{proof}
الحالتان $\{p, q\} = \{1, \infty\}$ مباشرتان ($\abs{fg} \leq
\norm g_\infty\abs f$ في كل مكان تقريبًا). وليكن
$1 < p < \infty$؛ نعيِّر $\norm f_p = \norm g_q = 1$ (بالتجانس؛ فالنظائم المعدومة أو
اللانهائية بديهية). وتعطي متراجحة يونغ $ab \leq \frac{a^p}p +
\frac{b^q}q$ (حيث
$a, b \geq 0$؛ بتقعّر $\ln$، كما في \cref{pb:b3:banach:1}) نقطةً نقطة أن
$\abs{f g} \leq
\frac{\abs f^p}p + \frac{\abs g^q}q$؛ ونكامل: $\norm{fg}_1
\leq \frac1p + \frac1q = 1$. ويفرض التساوي تساويًا في متراجحة يونغ
في كل مكان تقريبًا، أي $\abs f^p = \abs g^q$ في كل مكان تقريبًا (بعد التعيير؛
وبالرجوع عنه، التناسب).
\end{proof}

\begin{theorem}[متراجحة مينكوفسكي]\label{thm:b3:lp:minkowski}
من أجل $1 \leq p \leq \infty$: $\norm{f + g}_p \leq \norm f_p +
\norm g_p$.\index{متراجحة مينكوفسكي}
\end{theorem}

\begin{proof}
الحالتان $p = 1, \infty$: بمتراجحة المثلث نقطةً نقطة أو في كل
مكان تقريبًا. ومن أجل $1 <
p < \infty$، نفترض أن $\norm{f+g}_p < \infty$ (وإلا استعملنا
$\abs{f+g}^p \leq 2^{p-1}(\abs f^p + \abs g^p)$، الناتجة عن تحدّب $t^p$، لنرى أن الطرف الأيسر منتهٍ حين
يكون الأيمن كذلك). عندئذٍ
\[
\norm{f{+}g}_p^p \leq \int\abs f\,\abs{f{+}g}^{p-1} +
\int\abs g\,\abs{f{+}g}^{p-1}
\leq \bigl(\norm f_p + \norm g_p\bigr)\,
\bigl\|\abs{f{+}g}^{p-1}\bigr\|_q
\]
بهولدر، و $\norm{\abs{f+g}^{p-1}}_q =
\norm{f+g}_p^{p/q}$ لأن $(p-1)q = p$؛ ثم نقسم على
$\norm{f+g}_p^{p/q}$ (إن لم يكن معدومًا؛ وإلا فالنتيجة بديهية) ونستعمل $p -
\frac pq = 1$.
\end{proof}

\section{التمام وما يصحبه}

\begin{theorem}[ريس--فيشر]\label{thm:b3:lp:complete}
من أجل $1 \leq p \leq \infty$، يكون $L^p(\mu)$ فضاء باناخ. وأكثر من ذلك، لكل
متتالية متقاربة في $L^p$ متتاليةٌ جزئية متقاربة \emph{في كل مكان
تقريبًا} (مع مهيمِن في $L^p$ في
حالة $p < \infty$).\index{مبرهنة ريس--فيشر}
\end{theorem}

\begin{proof}
حالة $p = \infty$: كل متتالية كوشية بالمعنى $\norm\cdot_\infty$ هي، خارج
مجموعة معدومة واحدة (أي اتحاد عدد قابل للعدّ منها)، كوشيةٌ بانتظام:
فتتقارب بانتظام خارجها؛ وانتهى. وليكن $p < \infty$. فحسب
\cref{exo:b3:complete:1}(b) يكفي جمع المتسلسلات المتقاربة تقاربًا مطلقًا:
ليكن $\sum\norm{f_k}_p = M < \infty$. نضع $G_n = \sum_{k \leq n}\abs{f_k}$ و $G = \sum_k\abs{f_k}$ (نقطةً نقطة في $[0,\infty]$):
فبمينكوفسكي $\norm{G_n}_p \leq
M$، وتعطي مبرهنة التقارب الرتيب ($G_n^p \nearrow G^p$) أن
$\int G^p \leq M^p$: أي $G < \infty$ في كل مكان تقريبًا، ومنه تتقارب
المتسلسلة $\sum f_k(x)$ تقاربًا مطلقًا من أجل $x$ في كل مكان
تقريبًا؛ ولنسمّ مجموعها $S(x)$ (بأي قيمة على المجموعة المعدومة).
عندئذٍ $\abs{S - \sum_{k\leq n}f_k}^p \leq (2G)^p
\in L^1$، وتعطي مبرهنة التقارب المهيمن أن $\norm{S - \sum_{k\leq n}f_k}_p \to 0$: أي إن
المتسلسلة تتقارب في $L^p$.

وأما عبارة المتتالية الجزئية: فإذا كان $f_n \to f$ في $L^p$،
اخترنا $n_k$ تحقق $\norm{f_{n_{k+1}} - f_{n_k}}_p \leq 2^{-k}$؛ فتقع المتسلسلة $\sum(f_{n_{k+1}} - f_{n_k})$ تحت الحجة
السابقة: إذ تتقارب تقاربًا مطلقًا في كل مكان تقريبًا، وتُهيمَن
بمقدار $G \in
L^p$، ومنه $f_{n_k} \to f_{n_1} + \sum(\cdots)$ في كل مكان تقريبًا، ولا بد أن تكون
هذه النهاية ممثلًا للدالة $f$ (إذ كلتاهما نهاية في $L^p$). وأما
المهيمِن: فهو $\abs{f_{n_k}} \leq \abs{f_{n_1}} + G$.
\end{proof}

\begin{remark}\label{rem:b3:lp:modes}
لا يستلزم التقارب في $L^p$ التقاربَ في كل مكان تقريبًا (كما في
متتالية \emph{الآلة الكاتبة}، \cref{exo:b3:lp:3})، ولا العكس (كما في
النتوءات الهاربة): فالنمطان لا يرتبطان إلا عبر المتتاليات
الجزئية والهيمنة. واستحضار الأمثلة المضادة في \cref{exo:b3:lp:3} خير
لقاح ضد الخلط.
\end{remark}

\section{مبرهنات الكثافة}

\begin{theorem}\label{thm:b3:lp:density}
ليكن $1 \leq p < \infty$.
\begin{enumerate}
  \item الدوال البسيطة (ذات الحوامل المنتهية \omterm{def:b3:measure:measure}{القياس}) كثيفةٌ
        في $L^p(\mu)$.
  \item في $L^p(\R^d)$، تكون الدوال المتصلة ذات الحامل
        المتراص $\mathcal C_c(\R^d)$ كثيفة.
  \item الانسحاب متصل على $L^p(\R^d)$: فبكتابة
        $\tau_hf = f(\cdot - h)$، يكون $\norm{\tau_hf - f}_p \to 0$
        حين $h \to 0$.
\end{enumerate}
ولا تصح أي واحدة من الثلاث من أجل $p = \infty$.
\end{theorem}

\begin{proof}
(1) من أجل $f \geq 0$: تحقق الدوال الثنائية $s_n \nearrow f$ من
\cref{thm:b3:lebesgue:approximation} أن $\abs{f - s_n}^p
\leq f^p \in L^1$: بالتقارب المهيمن. (وكل $s_n \leq f$ ينتمي
إلى $L^p$، ومجموعات مستوياته منتهية \omterm{def:b3:measure:measure}{القياس} حيث تكون القيمة
موجبة: $\mu(s_n \geq c) \leq c^{-p}\int f^p$.) ثم نشطر $f$ العامة إلى أربعة أجزاء غير سالبة.

(2) حسب (1) يكفي تقريب $\mathbf 1_A$، حيث $A$ بوريلية تحقق
$\lambda_d(A) < \infty$. ويعطي الانتظام (بالبرهان نفسه كما في \cref{thm:b3:measure:regularity}) متراصةً
$K \subseteq A
\subseteq U$ ومفتوحةً تحققان $\lambda_d(U\setminus K) <
\varepsilon$؛ ودالة أوريسون
\[
\varphi(x) = \frac{d(x, \R^d\setminus U)}{d(x, \R^d\setminus U)
+ d(x, K)}
\]
متصلة، وتساوي $1$ على $K$ و $0$ خارج $U$، ويمكن أخذها ذات حامل
متراص (بتقليص $U$ إلى مفتوحة محدودة أولًا). عندئذٍ
$\norm{\mathbf 1_A - \varphi}_p^p \leq \lambda_d(U\setminus K)
< \varepsilon$.

(3) من أجل $g \in \mathcal C_c$: يعطي الاتصال المنتظم أن $\norm{\tau_hg - g}_\infty \to 0$، مع حوامل في
متراصة ثابتة من أجل $\abs h \leq 1$: ومنه $\norm{\tau_hg - g}_p \to 0$. وأما من أجل
$f$ عامة: فنختار $g \in \mathcal C_c$ تحقق $\norm{f - g}_p <
\varepsilon$؛ عندئذٍ $\norm{\tau_hf - f}_p \leq 2\norm{f - g}_p +
\norm{\tau_hg - g}_p$ (بصمود
المعيار بالانسحاب).

وأما من أجل $p = \infty$: فالتقريب المنتظم للدالة $\mathbf
1_{\intoo0\infty}$ بدوال
متصلة مستحيل (بسبب القفزة)، ويكون $\norm{\tau_h\mathbf 1_{\intoo0\infty} - \mathbf
1_{\intoo0\infty}}_\infty = 1$ من أجل $h \neq 0$.
\end{proof}

\section{الالتفاف والتنظيم}

\begin{theorem}[متراجحة يونغ]\label{thm:b3:lp:young}
ليكن $1 \leq p \leq \infty$، $f \in L^1(\R^d)$، $g \in
L^p(\R^d)$. عندئذٍ تكون $f * g$ معرَّفة في كل
مكان تقريبًا، وتنتمي إلى $L^p$،
و\index{متراجحة يونغ}
\[
\norm{f * g}_p \leq \norm f_1\,\norm g_p .
\]
\end{theorem}

\begin{proof}
حالة $p = \infty$: بالحد المباشر. وحالة $p = 1$: هي
\cref{thm:b3:product:convolution}. وليكن $1 < p < \infty$ وليكن $q$ المرافق. نشطر
$\abs{f(y)} = \abs{f(y)}^{1/q}\cdot
\abs{f(y)}^{1/p}$ ونطبّق هولدر:
\[
\int\abs{f(y)}\,\abs{g(x{-}y)}\,\dd y
\leq \Bigl(\int\abs f\Bigr)^{1/q}
\Bigl(\int\abs{f(y)}\,\abs{g(x - y)}^p\,\dd y\Bigr)^{1/p} .
\]
ثم نرفع إلى القوة $p$ ونكامل في $x$؛ وتعطي تونيلي على العامل
الثاني أن $\norm f_1^{p/q}\cdot\norm f_1\norm g_p^p$، أي $\norm{f*g}_p^p \leq \norm f_1^{1 + p/q}\norm g_p^p =
(\norm f_1\norm g_p)^p$ --- ويبرّر انتهاءُ تكامل تونيلي
التقاربَ المطلق في كل مكان تقريبًا كما في \cref{thm:b3:product:convolution}.
\end{proof}

\begin{definition}[المليِّفات]\label{def:b3:lp:mollifier}
الدالة\index{مليِّف}\index{دالة نتوء}
\[
\rho(x) = \begin{cases}
c\,\exp\Bigl(-\dfrac{1}{1 - \norm x^2}\Bigr) & \norm x < 1,\\
0 & \norm x \geq 1,
\end{cases}
\]
حيث يعيِّر $c$ الشرط $\int\rho = 1$، هي $\mathcal C^\infty$ على $\R^d$: إذ
المهم أن $t \mapsto \eu^{-1/t}\mathbf
1_{t>0}$ دالة $\mathcal C^\infty$ على $\R$، وأن جميع مشتقاتها عند
$0^+$ معدومة (فكل مشتقة من الشكل $P(1/t)\eu^{-1/t}$ من أجل كثير حدود $P$،
وهو يؤول إلى $0$؛ بالتراجع). ومن أجل $\varepsilon > 0$ نضع $\rho_\varepsilon(x) =
\varepsilon^{-d}\rho(x/\varepsilon)$: فحاملها
في $\bar B(0,
\varepsilon)$، وتكاملها لا يزال $1$.
\end{definition}

\begin{theorem}[التنظيم]\label{thm:b3:lp:regularization}
ليكن $1 \leq p < \infty$ و $f \in L^p(\R^d)$. عندئذٍ:
\begin{enumerate}
  \item $f * \rho_\varepsilon \in \mathcal C^\infty(\R^d)$،
        مع $\partial^\alpha(f * \rho_\varepsilon) = f *
        \partial^\alpha\rho_\varepsilon$؛
  \item $\norm{f * \rho_\varepsilon - f}_p \to 0$ حين
        $\varepsilon \to 0$؛
  \item ومن ثَمّ تكون $\mathcal C^\infty_c(\R^d)$ كثيفة في
        $L^p(\R^d)$.
\end{enumerate}
\end{theorem}

\begin{proof}
(1) بالاشتقاق تحت علامة التكامل (\cref{thm:b3:lebesgue:paramdiff}) في $x$: فمن أجل
$x$ في كرة $B$، يكون $\abs{\partial_{x_i}\rho_\varepsilon(x - y)} \leq
C_\varepsilon\,\mathbf 1_{K}(y)$ حيث $K$ متراصة (أي النقاط $y$
البعيدة عن $B$ بأقل من $\varepsilon$)، ويكون $\abs f\,\mathbf 1_K \in L^1$ (بهولدر
مقابل $\mathbf 1_K$): فتنطبق المبرهنة؛ ونكرّر من أجل المشتقات
العليا.

(2) بما أن $\int\rho_\varepsilon = 1$:
\[
(f * \rho_\varepsilon)(x) - f(x)
= \int \bigl(f(x - y) - f(x)\bigr)\rho_\varepsilon(y)\,\dd y ,
\]
ومتراجحة مينكوفسكي \emph{التكاملية} --- أو مباشرةً: هولدر أو
ينسن مع \omterm{def:b3:measure:measure}{القياس} الاحتمالي $\rho_\varepsilon\dd y$ ثم تونيلي ---
\[
\norm{f*\rho_\varepsilon - f}_p^p
\leq \int\Bigl(\int\abs{f(x-y) -
f(x)}^p\dd x\Bigr)\rho_\varepsilon(y)\,\dd y
= \int \norm{\tau_yf - f}_p^p\;\rho_\varepsilon(y)\,\dd y
\]
(والخطوة الوسطى: نطبّق متراجحة ينسن، \cref{exo:b3:lp:10}، على التكامل
الداخلي في $y$، ثم تونيلي). وحامل المقدار المكامَل في $\norm y \leq \varepsilon$،
وهو يؤول إلى $0$ هناك بانتظام حين $\varepsilon \to 0$ (\cref{thm:b3:lp:density}(3)): ومنه
يؤول التعبير كله إلى $0$.

(3) نقرّب $f$ بدالة $g \in \mathcal C_c$ (\cref{thm:b3:lp:density}(2))، ثم نقرّب $g$
بالمقدار $g *
\rho_\varepsilon \in \mathcal C_c^\infty$ (وحامله متراص: لأنه مجموع الحاملين).
\end{proof}

\begin{example}[تمليف $\abs x$، مع المعدلات]
\label{ex:b3:lp:mollifyexample}
لنأخذ $f(x) = \abs x$ على $\R$ (وهي في $L^1$ محليًّا؛ فتنطبق
المبرهنة على كل نافذة محدودة) ومليِّفًا متناظرًا $\rho_\varepsilon$. عندئذٍ
\[
f_\varepsilon(x) = (f * \rho_\varepsilon)(x)
= \int\abs{x - y}\,\rho_\varepsilon(y)\,\dd y
\]
تكون $\mathcal C^\infty$؛ وبعيدًا عن الانكسار لا يحدث شيء: فمن أجل $\abs x \geq \varepsilon$،
يكون $\abs{x - y}$ خطيًّا في $x$ على حامل $\rho_\varepsilon$، ومنه $f_\varepsilon(x) =
\abs x$
\emph{بالضبط} (إذ يلغي التناظر التصحيح). وبجوار $0$، يكلّف
التنعيم بالضبط
\[
0 \leq f_\varepsilon(0) = \int\abs
y\,\rho_\varepsilon(y)\,\dd y \leq \varepsilon,
\qquad
\norm{f_\varepsilon - f}_\infty \leq \varepsilon :
\]
أي إن خطأ التقريب محصور في جوار السويّة $\varepsilon$ للشذوذ،
وهو من حجمه. وفي الوقت نفسه $f_\varepsilon'' \geq 0$ في كل مكان (لأن $f$ محدّبة،
والالتفاف مع $\rho_\varepsilon
\geq 0$ يحفظ التحدّب)، مع $\int f_\varepsilon'' =
f_\varepsilon'(\infty) - f_\varepsilon'(-\infty) = 2$: فالمشتقة الثانية
نتوءٌ كتلته $2$ مضغوطٌ في عرض $O(\varepsilon)$، ومنه $\norm{f_\varepsilon''}_\infty \gtrsim
\varepsilon^{-1}$. فالتنعيم مقايضة: خطأ
منتظم من رتبة $O(\varepsilon)$ مقابل انفجار المشتقة $O(\varepsilon^{-1})$ ---
وهو سعر الصرف الدقيق الذي يصوغه التحليل الكمّي (متراجحات
الاستيفاء، ودائرة الأفكار في \cref{pb:b3:lp:1}).
\end{example}

\begin{corollary}[المبرهنة المساعدة الأساسية في حساب
التغيّرات]\label{cor:b3:lp:fundlemma}
لتكن $f \in L^1_{\mathrm{loc}}(\R^d)$ (\omterm{def:b3:lebesgue:l1}{قابلة للمكاملة} على المتراصات) وتحقق $\int f\varphi = 0$ من أجل
كل $\varphi \in \mathcal
C^\infty_c(\R^d)$. عندئذٍ $f = 0$ في كل مكان تقريبًا.\index{المبرهنة
المساعدة الأساسية في حساب التغيّرات}
\end{corollary}

\begin{proof}
نثبّت كرة $B = B(0, R)$ ولتكن $g = f\mathbf 1_{B(0, R+1)}
\in L^1$. من أجل $x \in B$ ومن أجل
$\varepsilon < 1$: $(g *
\rho_\varepsilon)(x) = \int f(y)\rho_\varepsilon(x - y)\dd y =
0$، ودالة الاختبار هي $y \mapsto \rho_\varepsilon(x-y)
\in \mathcal C_c^\infty$. لكن $g * \rho_\varepsilon \to g$ في
$L^1$ (\cref{thm:b3:lp:regularization}): ومنه $g = 0$ في كل مكان تقريبًا على $B$؛ ثم
نستنفد $\R^d$.
\end{proof}

\section{جغرافيا \texorpdfstring{$L^p$}{Lp}}

\begin{proposition}\label{prop:b3:lp:inclusions}
(a) إذا كان $\mu(X) < \infty$ و $1 \leq p \leq q \leq \infty$، فإن $L^q \subseteq L^p$ مع $\norm f_p \leq
\mu(X)^{\frac1p - \frac1q}\,\norm f_q$.
(b) وعلى $\R^d$ (ذي \omterm{def:b3:measure:measure}{القياس} اللانهائي) لا توجد أي احتواءات:
فمن أجل $p \neq q$ توجد دوال في $L^p\setminus L^q$.
(c) (الاستيفاء) إذا كان $p < r < q$ وعُرِّف $\alpha \in \intoo01$ بالعلاقة
$\frac1r = \frac\alpha p + \frac{1 - \alpha}q$، فإن
\[
\norm f_r \leq \norm f_p^{\alpha}\,\norm f_q^{1 - \alpha} ;
\]
وعلى وجه الخصوص $L^p \cap L^q \subseteq L^r$.
\end{proposition}

\begin{proof}
(a) بهولدر مع الأُسّين $\frac qp$ ومرافقه:
$\int\abs f^p\cdot 1 \leq \norm{\abs f^p}_{q/p}\,\norm
1_{(q/p)'} = \norm f_q^p\,\mu(X)^{1 - p/q}$ (وحالة $q = \infty$ مباشرة). (b) بجوار $0$ وبجوار
$\infty$، تعاير القوى $x^{-\alpha}$ الوضع: \cref{exo:b3:lp:2}.
(c) نكتب $\abs f^r = \abs
f^{r\alpha}\abs f^{r(1-\alpha)}$ ونطبّق هولدر مع الزوج المترافق $\frac p{r\alpha}$
و $\frac q{r(1-\alpha)}$ (وهما مترافقان بحكم تعريف $\alpha$ بالضبط):
$\int\abs f^r \leq \norm f_p^{r\alpha}\norm f_q^{r(1 -
\alpha)}$.
\end{proof}

\begin{method}\label{met:b3:lp:toolkit}
عدّة $L^p$، كما تُستعمل في كل ما يأتي: لتبرهن على متطابقة أو
متراجحة من أجل كل $f \in L^p$ --- برهن عليها على صنف كثيف
($\mathcal C_c^\infty$ عبر \cref{thm:b3:lp:regularization}) ومدّدها بالاتصال (\cref{thm:b3:complete:extension}، إذ
الطرفان متصلان بالمعنى $L^p$)؛ ولتبرهن على أن $f = 0$، اختبر
مقابل $\mathcal
C_c^\infty$ (\cref{cor:b3:lp:fundlemma})؛ ولتكسب نعومة، التفّ؛ ولتقايض
الأُسّ، استعمل هولدر والاستيفاء. ونظرية فورييه في \cref{ch:b3:fouriertransform}
تطبيقٌ واحد طويل لهذه الطريقة.
\end{method}

\section{تمارين}

\begin{exercise}[$\star$]\label{exo:b3:lp:1}
(a) صُغ متراجحة كوشي--شوارتز في $L^2(\mu)$ وبرهن عليها بوصفها
حالة $p = q = 2$ من هولدر.
(b) على فضاء \emph{احتمالي}، برهن على أن $p \mapsto \norm f_p$ متزايدة
بالمعنى الواسع.
(c) متى تكون هولدر تساويًا من أجل $p = 1$ و $q = \infty$؟
\end{exercise}

\begin{exercise}[$\star$]\label{exo:b3:lp:2}
من أجل أي $p \in \intco1\infty$ تنتمي الدوال التالية إلى $L^p$؟
\[
x^{-1/2}\mathbf 1_{\intoo01},\qquad
x^{-1/2}\mathbf 1_{\intoo1\infty},\qquad
\frac{1}{x^{1/2}(1 + \abs{\ln x})}\ \text{على} \intoo01,
\qquad
\frac1{1 + \abs x}\ \text{على} \R .
\]
واخلص إلى: أن $p$ الصغير أيسر على $\intoo01$، وأن $p$ الكبير
أيسر على $\intoo1\infty$، وأنه لا يحتوي أي $L^p$ آخرَ على $\R$.
\end{exercise}

\begin{exercise}[$\star\star$]\label{exo:b3:lp:3}
(الآلة الكاتبة) عدّد الفترات الثنائية $I_1 =
\intcc01$، $I_2 = \intcc0{\frac12}$،
$I_3 = \intcc{\frac12}1$، $I_4 = \intcc0{\frac14}$، \dots\ ولتكن $f_n = \mathbf
1_{I_n}$.
(a) برهن على أن $f_n \to 0$ في كل $L^p(\intcc01)$
حيث $p < \infty$، لكن $(f_n(x))$ تتباعد من أجل \emph{كل}
$x \in \intcc01$.
(b) أبرز المتتالية الجزئية المتقاربة في كل مكان تقريبًا التي
تعد بها \cref{thm:b3:lp:complete}.
(c) وبالعكس، أعطِ متتالية تتقارب في كل مكان تقريبًا دون أن
تتقارب في $L^1$، وأخرى تتقارب في $L^1$ دون أن تتقارب في أي
$L^p$ حيث $p > 1$.
\end{exercise}

\begin{exercise}[$\star\star$]\label{exo:b3:lp:4}
ليكن $\mu(X) < \infty$ و $f \in L^\infty(\mu)$، مع $f \neq 0$. برهن على أن $\norm f_p \to \norm f_\infty$ حين
$p \to \infty$. \emph{(الحدّ الأعلى من (a) في \cref{prop:b3:lp:inclusions}؛
والحدّ الأدنى بالمكاملة على $\{\abs f > \norm f_\infty -
\varepsilon\}$، وهي ذات \omterm{def:b3:measure:measure}{قياس} موجب.)}
\end{exercise}

\begin{exercise}[$\star\star$]\label{exo:b3:lp:5}
(a) أين بالضبط يستعمل برهان \cref{thm:b3:lp:density}(3) الفرضيةَ
$p < \infty$؟
(b) برهن على أن $f \in L^\infty(\R)$ تحقق $\norm{\tau_hf -
f}_\infty \to 0$ إذا وفقط إذا كان للدالة $f$
ممثّل متصل بانتظام.
\end{exercise}

\begin{exercise}[$\star\star$]\label{exo:b3:lp:6}
ليكن $p, q$ مترافقين، وليكن $f \in L^p(\R^d)$ و $g \in
L^q(\R^d)$. برهن على أن
$f * g$ معرَّفة \emph{في كل نقطة}، ومحدودة، مع $\norm{f*g}_\infty \leq \norm f_p\norm g_q$،
و\emph{متصلة بانتظام}. \emph{(باتصال الانسحاب في $L^p$؛ وعالج
$p \in \{1, \infty\}$ على حدة --- إذ من أجل $p =
\infty$ يُستعمل اتصال الانسحاب على
العامل الذي في $L^1$.)}
\end{exercise}

\begin{exercise}[$\star\star$]\label{exo:b3:lp:7}
لتكن $f \in L^1_{\mathrm{loc}}(\intoo ab)$ تحقق $\int f\varphi' = 0$ من أجل كل $\varphi \in \mathcal
C^\infty_c(\intoo ab)$. برهن على أن $f$
تساوي ثابتًا في كل مكان تقريبًا. \emph{(ثبّت $\chi \in \mathcal C_c^\infty$ تحقق
$\int\chi = 1$؛ فكل $\psi \in \mathcal C_c^\infty$ تحقق $\int\psi = 0$ هي $\varphi'$
ما؛ ثم اكتب دالة اختبار عامة على الشكل $\psi + (\int\psi)\chi$ وطبّق
\cref{cor:b3:lp:fundlemma} على $f - c$ حيث $c = \int
f\chi$.)}
\end{exercise}

\begin{exercise}[$\star\star\star$]\label{exo:b3:lp:8}
(أوريسون الناعمة) ليكن $K \subseteq U \subseteq \R^d$، مع $K$ متراصة و $U$ مفتوحة.
أنشئ $\varphi \in \mathcal
C^\infty_c(\R^d)$ تحقق $0 \leq \varphi \leq 1$، و $\varphi = 1$ على $K$، و $\operatorname{supp}\varphi \subseteq U$.
\emph{(بتمليف الدالة المميّزة \omterm{def:b3:topology:topology}{لجوار} السويّة $\delta$ للمجموعة
$K$، أي $K_\delta$، \omterm{def:b3:lp:mollifier}{بالمليِّف} $\rho_{\delta/2}$، من أجل $\delta$ صغير.)}
واستنتج عبارة تجزئة الوحدة من الصنف $\mathcal C^\infty$ من أجل متراصة
مغطّاة بعدد منتهٍ من المفتوحات.
\end{exercise}

\begin{exercise}[$\star\star$]\label{exo:b3:lp:9}
باستعمال الاستيفاء (\cref{prop:b3:lp:inclusions}(c)):
(a) برهن على أن $L^1(\R)\cap L^\infty(\R) \subseteq L^p(\R)$ من أجل كل $p$، مع $\norm f_p \leq \norm f_1^{1/p}\norm
f_\infty^{1 - 1/p}$؛
(b) برهن على أن $f \mapsto \norm f_p$، من أجل $f$ مثبَّتة، محدّبةٌ لوغاريتميًا
في $\frac1p$، وأعطِ مثالًا يكون فيه $f \in L^p$ من أجل $p$
في فترة معطاة $(p_0, p_1)$ بالضبط.
\end{exercise}

\begin{exercise}[$\star\star$]\label{exo:b3:lp:10}
(ينسن) ليكن $\mu$ قياسًا \emph{احتماليًّا}، ولتكن $f \in
L^1(\mu)$
حقيقية، ولتكن $\Phi \colon \R \to \R$ محدّبة. برهن على
\[
\Phi\Bigl(\int f\,\dd\mu\Bigr) \leq \int \Phi\circ
f\,\dd\mu
\]
\emph{(بمستقيم الإسناد للدالة $\Phi$ عند النقطة $m = \int f$)}.
واستنتج متراجحة الوسطين الحسابي والهندسي، ورتابةَ المقدار
$p \mapsto \norm f_p$ في \cref{exo:b3:lp:1}(b).\index{متراجحة ينسن}
\end{exercise}

\begin{exercise}[$\star\star\star$]\label{exo:b3:lp:11}
(متراجحة يونغ للالتفاف) ليكن $1 \leq p, q, r \leq
\infty$ مع $\frac1p + \frac1q = 1 + \frac1r$، وليكن $f \in
L^p(\R^d)$
و $g \in L^q(\R^d)$.
(a) برهن على $\norm{f * g}_r \leq \norm f_p\,\norm g_q$.
\emph{(اكتب، من أجل الأُسّين المترافقين المستخرجين من $p, q,
r$،
\[
\abs{f(y)g(x-y)} = \bigl(\abs f^p\abs g^q\bigr)^{1/r}
\cdot\abs f^{\,p(1/p - 1/r)}\cdot\abs g^{\,q(1/q - 1/r)},
\]
وطبّق متراجحة هولدر ذات العوامل الثلاثة بالأُسس $r$ و $\frac{pr}{r - p}$
و $\frac{qr}{r - q}$؛ ثم كامل في $x$ بتونيلي.)}
(b) تحقق من الحالات الخاصة الثلاث المعروفة سلفًا: $r =
\infty$
(هولدر، \cref{exo:b3:lp:6})؛ و $q = 1$ (أي استقرار $L^p$ بالالتفاف مع
نواة \omterm{def:b3:lebesgue:l1}{قابلة للمكاملة})؛ و $p = q = 1$ (أي إن $L^1$ جبرٌ التفافي،
\cref{thm:b3:product:convolution}).
(c) ولماذا لا توجد متراجحة مع $\frac1p + \frac1q < 1 +
\frac1r$؟ \emph{(اختبر على
التمدّدات $f_\lambda(x) =
f(\lambda x)$ وقارن تحجيم الطرفين.)}
\end{exercise}

\begin{exercise}[$\star\star$]\label{exo:b3:lp:12}
(حالات التساوي) (a) في متراجحة هولدر $\int\abs{fg} \leq \norm f_p\norm g_q$ (حيث
$1 < p < \infty$)، برهن على أن التساوي يقع إذا وفقط إذا كان
$\abs f^p$ و $\abs g^q$ متناسبين في كل مكان تقريبًا.
\emph{(بتتبّع حالة التساوي في متراجحة يونغ $ab \leq \frac{a^p}p + \frac{b^q}q$، وهي
$a^p
= b^q$.)}
(b) وفي متراجحة مينكوفسكي $\norm{f + g}_p \leq \norm f_p
+ \norm g_p$ (حيث $1 < p < \infty$)، برهن
على أن التساوي مع $f,
g \neq 0$ يفرض $g = cf$ في كل مكان تقريبًا مع
$c > 0$.
(c) وقابل ذلك بالحالتين $p = 1$ و $p = \infty$: صف حالات
التساوي (الأوسع بكثير) هناك، على أمثلة.
\end{exercise}

\section{مسألة: متراجحة هاردي}

\begin{problem}[{مسألة نهاية الأسبوع --- متراجحة هاردي وثابتها
الأمثل}]\label{pb:b3:lp:1}
من أجل $f \in L^p(\intoo0{+\infty})$ حيث $1 < p < \infty$، نعرّف \emph{مؤثر هاردي}
\[
(Hf)(x) = \frac1x\int_0^x f(t)\,\dd t .
\]
وتؤكد متراجحة هاردي (1920) أن\index{متراجحة هاردي}
\[
\norm{Hf}_p \;\leq\; \frac{p}{p-1}\,\norm f_p ,
\]
وأن الثابت $\frac p{p-1}$ أمثلي وغير مبلوغ. وتبرهن هذه المسألة
على كل ذلك، ثم تمدّده إلى المتسلسلات.

\medskip
\noindent\textbf{الجزء الأول --- المتراجحة.} نفترض أولًا أن
$f
\geq 0$ متصلة ذات حامل متراص في $\intoo0{+\infty}$، ولتكن $F(x) = \int_0^xf$.
\begin{enumerate}
  \item برهن على أن $Hf \in L^p$: فبجوار $0$ تنعدم $F$ على
        \omterm{def:b3:topology:topology}{جوار} للصفر؛ وبجوار $\infty$ تكون $F$ محدودة، ومنه
        $(Hf)(x) = O(1/x)$، وينتمي $x \mapsto \frac1x$ إلى $L^p(\intoo1{+\infty})$ من أجل $p > 1$.
  \item كامل بالتجزئة لتبرهن على
        \[
        \int_0^\infty \Bigl(\frac Fx\Bigr)^p\dd x
        = \frac{p}{p-1}\int_0^\infty\Bigl(\frac
        Fx\Bigr)^{p-1}f(x)\,\dd x .
        \]
        \emph{(اشتقّ $x^{1-p}F^p$؛ وتنعدم الحدود الحدّية ---
        وبرّر ذلك عند الطرفين.)}
  \item طبّق هولدر على الطرف الأيمن واستنتج $\norm{Hf}_p \leq \frac p{p-1}\norm f_p$ من أجل
        هذه $f$.
  \item مدّد إلى $L^p$ كله: من أجل $f \geq 0$، أنشئ $f_n$
        متصلة ذات حامل متراص في $\intoo0{+\infty}$، مع $0 \leq f_n \nearrow f$ في كل مكان
        تقريبًا \emph{(ابتر، ثم قرّب رتيبًا --- وبرّر البناء)}؛
        عندئذٍ $Hf_n \nearrow Hf$ نقطةً نقطة (بمبرهنة التقارب الرتيب داخل
        المتوسط) وتنقل مبرهنة التقارب الرتيب المتراجحةَ إلى
        النهاية. وأما من أجل $f$ ذات إشارة أو عقدية، فاخلص
        بواسطة $\abs{Hf} \leq H\abs f$.
\end{enumerate}

\medskip
\noindent\textbf{الجزء الثاني --- الأمثلية.}
\begin{enumerate}[resume]
  \item من أجل $A > 1$ لتكن $f_A(t) = t^{-1/p}\,\mathbf
        1_{\intcc1A}(t)$. احسب $\norm{f_A}_p^p = \ln A$، ومن أجل
        $1 \leq x \leq A$،
        \[
        (Hf_A)(x) = \frac p{p-1}\;x^{-1/p}\,
        \bigl(1 - x^{-(1 - 1/p)}\bigr).
        \]
  \item استنتج $\liminf_{A\to\infty} \norm{Hf_A}_p/
        \norm{f_A}_p \geq \frac{p}{p-1}$، واخلص إلى أن الثابت أمثلي.
  \item برهن على أن التساوي $\norm{Hf}_p = \frac
        p{p-1}\norm f_p$ مع $f \neq 0$ مستحيل.
        \emph{(بتتبّع حالة التساوي في هولدر في السؤال 3: فهي
        تفرض سلوكًا من نمط $f = cx^{-1/p}$، وهو ليس في
        $L^p$.)}
\end{enumerate}

\medskip
\noindent\textbf{الجزء الثالث --- المتراجحة المتقطّعة.}
\begin{enumerate}[resume]
  \item من أجل $g \geq 0$ متناقصة بالمعنى الواسع على
        $(0,\infty)$ ومن أجل $a_n = g(n)$، قارن $\sum a_n^p$
        بالمقدار $\int g^p$، وقارن متوسطات $H$ تبعًا لذلك،
        لتستنتج من الجزء الأول \emph{متراجحة هاردي المتقطّعة}:
        من أجل $a_n \geq 0$،
        \[
        \sum_{n\geq1}\Bigl(\frac{a_1 + \dots +
        a_n}{n}\Bigr)^{p}
        \;\leq\;
        \Bigl(\frac{p}{p-1}\Bigr)^{p}\,\sum_{n\geq1}a_n^p
        \]
        --- وبرهن عليها أولًا من أجل $(a_n)$ المتناقصة بالمعنى
        الواسع عبر المقارنة أعلاه، ثم أرجع الحالة العامة إلى
        المتناقصة بإعادة الترتيب \emph{(وسلّم، مع تبرير من سطر
        واحد، بأن ترتيب $(a_n)$ تنازليًّا لا يمكن إلا أن يزيد
        الطرف الأيسر مع إبقاء الأيمن على حاله)}.
  \item استنتج: أنه إذا كان $\sum a_n^p < \infty$ فإن متوسطات تشيزارو
        للمتتالية $(a_n)$ تنتمي بدورها إلى $\ell^p$ --- وأعطِ
        مثالًا (حيث $p = 2$) يكون فيه $(a_n) \in \ell^2$ دون أن تكون
        $a_n$ قابلة للجمع، ومع ذلك تظل هاردي تتحكم في
        المتوسطات.
\end{enumerate}

\medskip
\noindent\textbf{الجزء الرابع --- خاتمة.}
\begin{enumerate}[resume]
  \item برهن على أن متراجحة هاردي تخفق من أجل $p = 1$: فمع
        $f = \mathbf 1_{\intcc01}$، احسب $Hf$ ولاحظ أن $Hf \notin L^1$. وأين
        ينكسر البرهان؟
\end{enumerate}

\medskip
\noindent\textbf{الجزء الخامس --- الدالة العظمى، ومبرهنة
لوبيغ في الاشتقاق.} تأخذ هاردي المتوسطات انطلاقًا من المبدأ؛
أما هاردي--ليتلوود فيأخذانها \emph{حول كل نقطة}. من أجل
$f \in L^1(\R)$ نعرّف
\[
Mf(x) = \sup_{r>0}\ \frac1{2r}\int_{x-r}^{x+r}\abs
f\,\dd\lambda .
\]
\begin{enumerate}[resume]
  \item (فيتالي، الصيغة المنتهية) لتكن $B_1, \dots, B_N$ فترات مفتوحة.
        برهن على وجود عائلة جزئية منفصلة $B_{i_1}, \dots, B_{i_k}$ تحقق $\bigcup_jB_j
        \subseteq \bigcup_l3B_{i_l}$،
        حيث يرمز $3B$ إلى الفترة ذات المركز نفسه وثلاثة أمثال
        الطول \emph{(بالجشع: اختر تكرارًا أطول فترة منفصلة عن
        المختارة سلفًا)}.
  \item (النمط الضعيف $(1,1)$) برهن على أنه من أجل كل $t > 0$،
        \[
        \lambda\bigl(\{Mf > t\}\bigr) \;\leq\;
        \frac3t\,\norm f_1 :
        \]
        فكل $x$ يحقق $Mf(x) > t$ يملك فترةً مركزيةً $B_x$
        تحقق $\int_{B_x}\abs f > t\,\lambda(B_x)$؛ خذ متراصة $K \subseteq \{Mf > t\}$ (بالانتظام الداخلي)،
        وغطّها بعدد منتهٍ من الفترات $B_x$، وطبّق السؤال 11،
        ثم استنفد.
  \item احسب $M\mathbf 1_{\intcc01}(x)$ من أجل $x > 1$ واستنتج أن $Mf \notin L^1$
        من أجل كل $f \neq 0$ (لأن $Mf(x) \geq \frac c{\abs x}$ عند اللانهاية): أي
        إن المتراجحة الضعيفة في السؤال 12 هي، عند $p =
        1$،
        أفضل عبارة ممكنة.
  \item (النمط القوي من أجل $p > 1$) من أجل $f \in L^p$: اشطر
        $f =
        f\,\mathbf 1_{\abs f > t/2} + f\,\mathbf 1_{\abs f
        \leq t/2}$، ولاحظ $Mf \leq M\bigl(f\mathbf 1_{\abs
        f > t/2}\bigr) + \frac t2$، واجمع السؤال 12 مع صيغة كعكة
        الطبقات (\cref{prop:b3:product:layercake}) وتونيلي لتبرهن على
        \[
        \norm{Mf}_p^p \;\leq\;
        \frac{6p\,2^{p-1}}{p-1}\,\norm f_p^p .
        \]
        (والانفجار حين $p \downarrow 1$ هو إخفاق السؤال 13،
        مكمَّمًا.)
  \item (مبرهنة لوبيغ في الاشتقاق) برهن على أنه من أجل
        $f \in
        L^1(\R)$،
        \[
        \frac1{2r}\int_{x-r}^{x+r}f\,\dd\lambda
        \;\longrightarrow\; f(x)
        \qquad (r \to 0)\quad\text{في كل مكان تقريبًا من أجل}x .
        \]
        \emph{(وهذا واضح من أجل $f$ متصلة. وأما في الحالة
        العامة فاكتب $f
        = g + h$، حيث $g$ متصلة ذات حامل متراص،
        مع
        $\norm h_1 < \varepsilon$
        (\cref{thm:b3:lp:density})؛ فتقع المجموعة التي يتجاوز فيها $\limsup_{r\to0}$
        للتذبذب الممتوسط العددَ $\delta$ داخل $\{Mh > \delta/2\} \cup \{\abs h
        > \delta/2\}$، وقياسها
        $O(\varepsilon/\delta)$؛ ثم اجعل $\varepsilon \to 0$، ثم $\delta \to 0$ على امتداد
        متتالية.)}
  \item استنتج: (a) أن كل نقطة تقريبًا \emph{نقطة لوبيغ}
        للدالة $f$؛ (b) وأنه من أجل $f \in L^1$، تكون الدالة
        الأصلية $F(x) = \int_0^xf$ قابلة للاشتقاق في كل مكان تقريبًا مع
        $F' =
        f$ في كل مكان تقريبًا --- أي النصف التكاملي
        للمبرهنة الأساسية في التفاضل والتكامل في عالم لوبيغ،
        فتُغلق بذلك الدائرة التي فتحها السلّم (\cref{pb:b3:measure:1})،
        إذ بيّن أن النصف العكسي قد يخفق.
  \item (نقاط الكثافة) من أجل $A \subseteq \R$ قابلة \omterm{def:b3:measure:measure}{للقياس}،
        برهن على أن كل $x \in A$ تقريبًا يحقق
        $\frac{\lambda(A\cap\intcc{x-r}{x+r})}{2r} \to 1$:
        أي إن المجموعات القابلة \omterm{def:b3:measure:measure}{للقياس} ممتلئة محليًّا عند كل
        نقاطها تقريبًا. واشرح، في سطرين، كيف يعطي هذا برهانًا
        آخر على مبرهنة شتاينهاوس (\cref{exo:b3:measure:8}).
\end{enumerate}

\medskip
\noindent\textbf{الجزء السادس --- تنويعات على موضوع
المتوسطات.}
\begin{enumerate}[resume]
  \item (هاردي الموزونة) من أجل $\alpha < p - 1$ ومن أجل
        $f \geq 0$، برهن على
        \[
        \int_0^\infty\Bigl(\frac{F(x)}x\Bigr)^{p}
        x^{\alpha}\,\dd x \;\leq\;
        \Bigl(\frac{p}{p - 1 - \alpha}\Bigr)^{p}
        \int_0^\infty f(x)^p\,x^{\alpha}\,\dd x
        \]
        بالمكاملة بالتجزئة نفسها، وتحقق من أن الحالة الحدّية
        $\alpha = p - 1$ ممنوعة فعلًا (بتكييف المثال المضاد في
        السؤال 10).
  \item (المؤثر المرافق) ليكن $H^*f(x) =
        \int_x^{\infty}\frac{f(t)}t\,\dd t$. برهن على $\langle
        Hf, g\rangle = \langle f, H^*g\rangle$ من أجل
        $f, g$ غير سالبتين (بتونيلي)، وبرهن على $\norm{H^*f}_p \leq p\,\norm f_p$
        \emph{(مباشرةً بالتجزئة، أو من هاردي على الأُسّ
        المرافق بالثنوية --- وانتبه إلى أي أُسّ يلتقط أي
        ثابت)}.
  \item (متراجحة من نمط هيلبرت) استنتج أنه من أجل $f \in L^p$
        و $g \in L^q$ غير سالبتين:
        \[
        \int_0^\infty\!\!\int_0^\infty
        \frac{f(x)\,g(y)}{\max(x,y)}\,\dd x\,\dd y
        \;\leq\; (p + q)\,\norm f_p\,\norm g_q
        \]
        \emph{(بالشطر على امتداد $y < x$ و $y \geq x$: فكل
        نصف اقترانٌ لإحدى الدالتين مع تحويل هاردي للأخرى)}.
  \item (الأمثلية، الحالة المتقطّعة) برهن على أن الثابت
        $\bigl(\frac p{p-1}\bigr)^p$ في السؤال 8 أمثلي أيضًا: اختبر على $a_n = n^{-1/p}\,\mathbf 1_{n \leq
        N}$،
        وقارن الطرفين بالتكاملات، ثم اجعل $N
        \to \infty$ (أي مرآة
        الجزء الثاني في الحالة المتقطّعة).
  \item (تركيب) ظهرت في هذه المسألة ثلاثة مؤثرات متوسطات:
        مؤثر هاردي $H$، ومتوسط تشيزارو المتقطّع، والمؤثر
        الأعظم $M$. صُغ في سطر واحد لكلٍّ منها ما يقوله كونه
        محدودًا، ولاحظ أن الثلاثة تخفق جميعًا عند $p = 1$
        بالضبط، واشرح لماذا هو الإخفاق نفسه ثلاث مرات (أي
        الذيل التوافقي $\frac1x$).
\end{enumerate}

\medskip
\noindent\textbf{الجزء السابع --- متراجحة كارلمان، وكم يبلغ
«الأمثل» من مضاء.}
\begin{enumerate}[resume]
  \item (متراجحة كارلمان) ليكن $a_n \geq 0$ مع $\sum
        a_n < \infty$. طبّق
        متراجحة هاردي المتقطّعة في السؤال 8 على $b_n = a_n^{1/p}$،
        واستعمل متراجحة الوسطين الحسابي والهندسي، ثم اجعل
        $p \to
        \infty$ (وبرهن على أن $p \mapsto
        \bigl(\frac p{p-1}\bigr)^p$ تتناقص إلى $\eu$)
        لتحصل على
        \[
        \sum_{n\geq1}\bigl(a_1a_2\cdots a_n\bigr)^{1/n}
        \;\leq\; \eu\,\sum_{n\geq1}a_n :
        \]
        أي إن المتوسطات الهندسية لمتتالية قابلة للجمع قابلةٌ
        للجمع، بكلفة لا تتجاوز $\eu$.
  \item (الثابت $\eu$ أمثلي) اختبر على $a_n =
        \frac1n\,\mathbf 1_{n\leq N}$: باستعمال حصر
        ستيرلنغ في \cref{pb:b3:product:1}، برهن على $(n!)^{-1/n} = \frac\eu n\bigl(1 +
        O\bigl(\frac{\ln n}n\bigr)\bigr)$، واستنتج أن
        طرفي كارلمان ينموان مثل $\eu\ln N$، واخلص إلى أنه لا
        يصلح أي ثابت أصغر من $\eu$. (ولاحظ النسق: فأمثِلات
        هاردي وكارلمان كلتاهما متتاليات من النمط التوافقي
        تكاد تنتمي إلى الفضاء دون أن تنتمي إليه.)
  \item (كم يبطؤ بلوغ «الأمثل»؟) خذ $p = 2$. من أجل $f = \mathbf 1_{\intcc01}$،
        احسب $\norm{Hf}_2/\norm f_2 = \sqrt2$، مقابل الحدّ $2$. ومن أجل الأمثِلات
        التقريبية $f_A$ في السؤال 5، برهن على المتطابقة
        المضبوطة
        \[
        \norm{Hf_A}_2^2 = 4\ln A - 8 + \frac{8}{\sqrt A},
        \qquad\text{ومنه}\qquad
        \frac{\norm{Hf_A}_2^2}{\norm{f_A}_2^2}
        = 4 - \frac{8 - 8A^{-1/2}}{\ln A} .
        \]
        وقيّمها عند $A = \eu^{10}$ (فتكون النسبة
        $\approx 1.790$) وعلّق: فالسوپريموم $2$ لا يُبلَغ إلا بسرعة
        $1/\ln A$ --- ومن ثَمّ قد يكون الثابت الأمثل شبه
        غير مرئي عدديًّا.
\end{enumerate}
\end{problem}
